\documentclass[12pt, a4paper]{article}
\usepackage{fontspec}
\usepackage{xeCJK}
\usepackage{geometry}
\usepackage{enumitem}
\usepackage{titlesec}
\usepackage{parskip}

\geometry{margin=2.5cm}
\setmainfont{Helvetica}
\setCJKmainfont{PingFang TC}

\title{約翰福音陪讀教材1}
\author{}
\date{}

\begin{document}

\maketitle

\section*{簡介}

約翰福音是四福音書中最後寫成的，風格與馬太、馬可、路加三卷（對觀福音）截然不同。
約翰以深刻的神學語言呈現耶穌的身份，核心主題是：耶穌是神的兒子，信他的人可以得永生。

\bigskip
\noindent\textbf{關鍵經文：}「神愛世人，甚至將他的獨生子賜給他們，叫一切信他的，不至滅亡，反得永生。」（約翰福音 3:16）

\section{第一章：道成肉身}

\subsection*{經文範圍}
約翰福音 1:1--18（序言）、1:19--51（施洗約翰與首批門徒）

\subsection*{默想問題}
\begin{enumerate}[leftmargin=*, label=\arabic*.]
  \item 「太初有道，道與神同在，道就是神」（1:1）——你如何理解「道」（Logos）這個概念？
  \item 約翰說他「不是那光，乃是要為光作見證」（1:8）。在你的生活中，你是如何為光作見證的？
  \item 耶穌對拿但業說：「你在無花果樹底下，我就看見你了」（1:48）。這句話讓你有什麼感受？
\end{enumerate}

\subsection*{生活應用}
想一想：這週有沒有一件事，你可以像施洗約翰一樣，指向耶穌而不是指向自己？

\section{第三章：重生與永生}

\subsection*{經文範圍}
約翰福音 3:1--21（尼哥底母）、3:22--36（約翰的謙卑）

\subsection*{默想問題}
\begin{enumerate}[leftmargin=*, label=\arabic*.]
  \item 尼哥底母是法利賽人、猶太人的官，他為什麼要夜間來見耶穌？這說明了什麼？
  \item 「人若不重生，就不能見神的國」（3:3）——重生在你的生命中是什麼經歷？
  \item 約翰說「那從地上來的，是屬地的」（3:31）。這如何幫助你分辨生命中不同的聲音？
\end{enumerate}

\subsection*{生活應用}
3:16 是整本聖經最廣為人知的經文之一。試著用自己的話，向一個從未聽過福音的朋友解釋這節經文。

\section{第四章：撒馬利亞婦人}

\subsection*{經文範圍}
約翰福音 4:1--42

\subsection*{默想問題}
\begin{enumerate}[leftmargin=*, label=\arabic*.]
  \item 耶穌突破了哪些文化與宗教的界限來與這位婦人說話？
  \item 「我所賜的水要在他裡頭成為泉源，直湧到永生」（4:14）——你生命中有這樣的泉源嗎？
  \item 婦人放下水罐回城報信（4:28），這個細節有什麼意義？
\end{enumerate}

\subsection*{生活應用}
這週留意身邊是否有「被社會邊緣化」的人，思考你可以如何主動靠近他們。

\bigskip
\hrule
\bigskip
\noindent\textit{本教材為個人陪讀使用，歡迎根據小組需要調整問題。}

\end{document}
